% Contenido científico íntegro normalizado del frontmatter y cortes
% Residencia material del ensamblaje sucesor de 102 capítulos.
% parte_i.tex:1--419,529--659 + parte_iii.tex:1--908.
% Se excluyen sólo la portadilla autónoma y sus órdenes de página.
\part{Genealogical radial correspondences on the enriched state: domain, codomain, and preservation of helical history}
This publication develops a genealogical theory of spiral
correspondences from the composition \(APP\to\TRIT\to\TPK\). The primary object
is not a curve presupposed in \(\R^2\), but a compatible history of the
enriched state, together with its arithmetic sheets, orientation, carry,
boundary, route, and memory. The curve appears at the end as an
Archimedean publication. Consequently, two identical figures in
\(\R^2\) may proceed from different histories, and a visible circle
may have intrinsically zero radial memory or be the shadow of
a helix whose memory was eliminated by the projection.

The work preserves and develops the exact results of the HMT corpus concerning
the mixed curvature of APP, the three quadratic regimes of TRIT, the
nonadic holonomy \(\Gamma_9\), phase return with memory advance,
the nonadic helix, the nonsplit dodecaphasic extension, and the Pell
spectral screw. On those results it constructs a new formalization:
the universal affine envelope, the prefix-enriched radial reader, its
naturality under truncation, its equivariance with \(\Gamma_9\), and its
power--logarithmic publication.

The construction does not alter the theory of masses or reopen the central electron.
The sole fixed point \((5,5)\) of the material atlas is used only as
the subsequent anchor of an interface; never as a retrospective selector of the
radial geometry.

\medskip
\begin{statusbox}[title={Probative stratification by domain and codomain}]
The exposition distinguishes previously established internal theorems, theorems
proved in this work, finite certifications, and physical realizations. An
internal algebraic result is asserted fully in its domain. Every
subsequent physical realization preserves its premises and does not intervene in the
selection of the HMT generator.
\end{statusbox}

\part{Construction and main results of the enriched radial reader}
A theory of spiral correspondences is constructed whose antecedent is
the enriched state of Triadic Modular Holography and whose output is a
family of curves in \(\R^2\) or \(\R^3\). APP provides a
two-dimensional support with two evaluations that are not simultaneously compatible,
\(\Sigma\) and \(\Pi\); their mixed composition has oriented discrete
curvature \(\pm1\). TRIT types the elliptic, parabolic, and
hyperbolic regimes without confusing them with the radial character or Frenet
curvature. TPK selects, transports, memorizes, updates, and returns the complete
state. The holonomy \(\Gamma_9\) returns the phase, but increments the
memory depth and transports the boundary.

The initially proposed reader into
\((p,a,\Lambda,C,m,r,\theta)\) is decisively corrected. The
coordinates \((r,\theta)\) belong to the Archimedean publication, and the
final pair \((\Lambda,C)\) does not preserve its prefixes. An internal
reader is therefore defined
\[
 \mathcal P^{\enr}_{\rad,N}:\Hist_N^{\enr}(\TPK)
 \longrightarrow \mathcal D^{\enr}_{\rad,N}
\]
which records, at each depth, the mode and orientation cocycles, the
local affine factor, prefix holonomy, and complete ledger. Its
naturality under truncation, extension to the inverse limit,
equivariance with \(\Gamma_9\), preservation of carry, boundary, and
memory, and compatibility with route inversion in the reversible
sector are proved.

After constructing the continuum, a radial character is applied and the
power--logarithmic coordinate is introduced
\[
 L_p(x)=\frac{x^p-1}{p},\qquad L_0(x)=\log x.
\]
Its composition law makes the APP sum--product articulation visible.
Uniform translation in \(L_p\) publishes, for
\(p=2,1,0,-1,-2\), the Fermat, Archimedean, logarithmic,
hyperbolic, and lituus spirals. General affine holonomy produces a broader
classification through \((p,[\Lambda,C])\). The spectral screw
\(V^{-1}\mathscr SV=(2+\sqrt3)\ii\mathscr S\) appears as an exact case:
its transverse projection is a logarithmic spiral, and its lift preserves
helical memory.

The publication concludes with the differential geometry of the published
curves, a restricted electromagnetic realization of the phase--position
helix, the interface with the central electron, and a system of
falsifiers preventing appearance, genealogy, and physical realization from being
conflated.

\bigskip
\noindent\textbf{Keywords:} APP; TRIT; Topological Phase Kernel;
nonadic holonomy; affine cocycle; spirals; helicoidal suspension; memory;
double projection; continuum; Dimensional Mechanics.

\part{Key results of the radial publication APP--TRIT--TPK}
\begin{enumerate}[label=\textbf{R\arabic*.},leftmargin=12mm]
\item Visible geometry is a nonfaithful publication of the enriched state:
equality in \(\R^2\) does not imply equality of genealogy.
\item The mixed curvature of APP distinguishes the sum--product order before
any real curve.
\item \(\tau\), \(p\) and the Frenet curvature are invariants of distinct levels
and cannot be identified.
\item The nonadic holonomy returns the phase, advances the memory and transports
the boundary; therefore the periodic angular projection admits a nonadic
helicoidal lift that preserves the depth of return.
\item Complete radial information requires preserving all factors and prefix holonomies; the final pair \((\Lambda,C)\) is insufficient.
\item The enriched radial reader is natural under truncation and
equivariant under \(\Gamma_9\).
\item The bidirectional involution produces conjugate ascending and descending routes in the reversible sector.
\item The family \(L_p\) publishes the five power--logarithmic classes as
outputs of one and the same law, without choosing them in advance.
\item The Pell spectral screw constitutes an exact realization of
rotation and self-scale whose projection is logarithmic.
\item A circumference may be a radially stationary state or the
shadow of a helix with nonzero ledger.
\item The electron \((5,5)\) remains the proven material center; the
spiral signature is added as a reader, not as a substitute for the atlas
\(81/56/13/324\).
\item The electromagnetic and material realizations are subsequent to the
HMT construction and preserve their own physical typing.
\end{enumerate}

\part{Typed notation and causal levels of the radial publication}
\begin{longtable}{>{\raggedright\arraybackslash}p{.23\textwidth}p{.67\textwidth}}
\toprule
Symbol & Meaning\\
\midrule
\endfirsthead
\toprule
Symbol & Meaning\\
\midrule
\endhead
\(X_9\) & Nonadic torus \((\Z/9\Z)^2\), the two-dimensional support of APP.\\
\(\Sigma,\Pi\) & Additive and multiplicative evaluations of the same support.\\
\(\tau\in\{-1,0,+1\}\) & Local TRIT type of transport: hyperbolic, parabolic, or elliptic according to the declared quadratic convention.\\
\(U_t\) & Typed transition \(\mathrm{Upd}_t\circ\mathrm{Tra}_t\circ\mathrm{Sel}_t\).\\
\(\Gamma_9\) & Holonomy of nine transitions over the enriched state.\\
\(p\) & Cocycle or radial degree after defining its internal extraction; it is not \(\tau\).\\
\(a\) & Local oriented character \(a=\tau_+-\tau_-\); it is distinct from the radial degree \(p=\tau_++\tau_-\).\\
\((\Lambda,C)\) & Radial affine endomorphism: linear part and translation.\\
\(H_j\) & Prefix affine holonomy up to depth \(j\).\\
\(L_p,\exp_p\) & Power-logarithmic coordinate and its inverse.\\
\(\kappa_{\mathrm{F}},\mathcal T_{\mathrm{F}}\) & Frenet curvature and torsion of the published curve.\\
\(\mathcal P^{\enr}_{\rad}\) & Radial internal genealogical reader.\\
\(\ev_{\arq}^{\mathrm{esp}}\) & Archimedean publication of the spiral firm.\\
\bottomrule
\end{longtable}

\begin{resultbox}[title={Causal composition from APP--TRIT--TPK to the radial Archimedean evaluation}]
\[
\APP\longrightarrow\TRIT\longrightarrow\TPK
\longrightarrow X_{\TPK}^{\enr}
\longrightarrow\mathcal C_{\mathrm{cont}}^{\mathrm{disc}}
\longrightarrow\mathcal P_{
  \rad}^{\enr}
\longrightarrow\ev_{\arq}^{\mathrm{esp}}
\longrightarrow\bigl(\R^2\ \text{or}\ \R^3\bigr).
\]
The curve is an output of the generative chain; it does not enter the selection
of the state or its coefficients.
\end{resultbox}
\part{Genealogical Foundation of Radial Geometry}

\chapter{Visible geometry as a publication of the enriched state}

\section{The inversion of the explanatory order}

Conventional geometry begins with a continuous set and defines in it
curves, parameters, and fields. The construction developed here reverses that
order. First a compatible history is produced on a finite support,
then its transport data are conserved, and only at the end is a real
coordinate published. The causal scheme is
\begin{equation}
 \APP\longrightarrow\TRIT\longrightarrow\TPK
 \longrightarrow X_{\TPK}^{\enr}
 \longrightarrow\mathcal C_{\mathrm{cont}}^{\mathrm{disc}}
 \longrightarrow\ev_{\arq}
 \longrightarrow\R^d.
 \label{espiral:eq:cadena-genealogica-espiral}
\end{equation}
It is not a matter of adorning a classical curve with discrete labels. The
history is the object; the curve is one of its readings.

\begin{definition}[Geometric Generation with Traceability]
To construction geometric is \emph{constructivo--generativa} when:
\begin{enumerate}
\item produces its object from declared operations on a domain
preceding real geometry;
\item preserves the order of operations and the coordinates needed to reconstruct it;
\item obtains the geometric appearance by means of a later map;
\item distinguishes equality of appearances from equivalence of histories.
\end{enumerate}
\end{definition}

This criterion is stronger than a parametrization. A parametrization
\(t\mapsto\gamma(t)\) presupposes the curve and changes its mode of traversal. A
genealogy \(x=(x_n)_{n\ge0}\), by contrast, specifies the finite truncations, the
admissible transitions, the arithmetic sheets, the carry, and the memory
that make the output \(\gamma\) possible.

\section{State, history and publication}

Let \(X_{\TPK}^{\enr}\) be the enriched state space. A finite
history of length \(N\) is
\[
 \gamma_N=(x_0\xrightarrow{e_1}x_1\xrightarrow{e_2}\cdots
 \xrightarrow{e_N}x_N),
\]
where every arrow is compatible with the typed TPK transition. We associate its ledger with each prefix
\[
 \Led_j(\gamma_N)=
 (R_j,Q_j,\kappa_j,\sigma_j,o_j,g_j,C_j,B_j,
  \lambda_j,m_j,s_j),
\]
which registers, according to the concrete residence, residue, quotient, carry,
leaf, orientation, phase, cylinder, boundary, route, memory, and survival.
The exact content may be enriched, but it cannot be reduced if the
eliminated coordinate intervenes in a continuation.

The real publication is a composition.
\[
 \Hist_N^{\enr}(\TPK)
 \xrightarrow{\ \mathcal P_N^{\enr}\ }
 \mathcal D_N^{\enr}
 \xrightarrow{\ \ev_N\ }
 Y,
\]
with \(Y=\R,\R^2\) or \(\R^3\). The first map preserves genealogy; the
second publishes. If \(\ev_N(\gamma)=\ev_N(\eta)\), it follows only that
both histories belong to the same evaluation fibre.

\begin{proposition}[No fidelity of the visible geometry]
If the evaluation map removes at least one ledger coordinate that
admits two compatible values, then there exist distinct histories with the
same geometric output.
\end{proposition}
\begin{proof}
Let \(c\) be an eliminated coordinate and let \(\gamma,\eta\) be two
compatible histories that coincide in all preserved coordinates and differ in
\(c\). By definition of the evaluation, \(\ev(\gamma)=\ev(\eta)\), whereas
\(\gamma\neq\eta\) because their ledgers differ. Therefore, \(\ev\) is not
injective.
\end{proof}

The proposition is not an insufficiency of real geometry. It expresses its
function: to contract a richer structure to a usable coordinate. The
central thesis consists precisely in not confusing that contraction with the
contracted object.

\section{Intrinsic circle and projected circle}

Consider the application
\[
 c_0(t)=(R\cos t,R\sin t),
\]
and the helix
\[
 h(t)=(R\cos t,R\sin t,bt).
\]
The projection \(\pi_{xy}(x,y,z)=(x,y)\) satisfies
\(\pi_{xy}\circ h=c_0\). The same circumference may represent:
\begin{enumerate}
\item a state with null radial increment and memory;
\item the shadow of a state with null radial increment and axial memory
\(b\neq0\).
\end{enumerate}
The visible form does not decide between both cases. The ledger is needed.

\begin{figure}[htbp]
\centering
\begin{subfigure}[t]{.46\textwidth}
\centering
\begin{tikzpicture}[scale=.82]
  \draw[->,HMTGray] (-2.2,0)--(2.2,0);
  \draw[->,HMTGray] (0,-2.2)--(0,2.2);
  \draw[HMTBlue,very thick] (0,0) circle (1.6);
  \fill[HMTNavy] (1.6,0) circle (1.7pt);
  \node[below,HMTGray,font=\small\sffamily] at (0,-2.25) {null visible memory};
\end{tikzpicture}
\caption{Circle as a planar object.}
\end{subfigure}\hfill
\begin{subfigure}[t]{.46\textwidth}
\centering
\begin{tikzpicture}[x={(.9cm,0cm)},y={(.32cm,.14cm)},z={(0cm,.72cm)},scale=.74]
  \draw[->,HMTGray] (0,0,0)--(0,0,6.8);
  \draw[HMTBlue,very thick,domain=0:720,samples=160,smooth,variable=\t]
    plot ({1.45*cos(\t)},{1.45*sin(\t)},{\t/110});
  \draw[HMTRed,dashed,thick] (0,0,0) ellipse (1.45 and .47);
  \node[right,HMTGray,font=\small\sffamily] at (0,0,6.8) {memory};
\end{tikzpicture}
\caption{The same shadow, with memory raised.}
\end{subfigure}
\caption{Visible geometry does not by itself reconstruct the genealogy.}
\label{espiral:fig:circulo-helice}
\end{figure}

\begin{resultbox}[title={Failure of fidelity of Archimedean evaluation on genealogical classes}]
A circle may have intrinsically null memory and may also
be the publication of a helical history. Equality of the curve does not
authorize identifying the states that produce it.
\end{resultbox}

\chapter{APP: two-dimensional support, arithmetic sheets and mixed curvature}

\section{The nonadic cell and the two-dimensional birth}

The visible cell
\[
 E_9=\{1,2,3,4,5,6,7,8,9\}
\]
simultaneously carries a phase and a chart of positive representatives. For
\(n\ge1\),
\[
 \rho_9(n)=1+((n-1)\bmod9),\qquad
 q_9(n)=\left\lfloor\frac{n-1}{9}\right\rfloor,
\]
so that
\[
 n=9q_9(n)+\rho_9(n).
\]
The Cartesian product of two cells constructs
\[
 X_9=(\Z/9\Z)^2.
\]
It is not a repeated line but a genuinely two-dimensional support. The
completion \(9\times9\) may be read as a core \(8\times8\) and a
gnomon of \(17\) cells:
\[
81=64+17.
\]

Two evaluations of the same support act on every \((i,j)\in X_9\):
\[
 \Sigma(i,j)=\operatorname{dr}(i+j),\qquad
 \Pi(i,j)=\operatorname{dr}(ij),
\]
where \(\operatorname{dr}\) publishes the nonadic representative and the
lift conserves the quotient. APP does not add two spaces; it conserves a
single space with two readings.

\section{Residue, quotient and carry}

The raised identities
\[
 i+j=R^+_{ij}+9q^+_{ij},\qquad
 ij=R^\times_{ij}+9q^\times_{ij}
\]
distinguish the visible output \(R\) from the lift \(q\). If only
the residue is preserved, distinct operations may collapse in the same
cell. The pair \((R,q)\) unfolds the history. Under composition, the
carry obeys a cocycle law; therefore it is not an auxiliary annotation,
but rather the datum that allows a long operation to be recomposed from its steps.

\begin{example}[Same residue, different history]
In the nonadic chart,
\[
 8+8=7+9\cdot1,\qquad 8\cdot2=7+9\cdot1.
\]
The two operations publish the same residue and the same local quotient,
but differ in leaf. The label \(\Sigma/\Pi\) remains necessary.
On longer routes, the order of those leaves modifies the composite carry and
the boundary reached.
\end{example}

\section{Operational Incompatibility of the Additive and Multiplicative Leaves}

For \(d\ge2\), let \(\mathscr A_d=E_9^d\). The conditional expectations
\(P_{\Sigma,d}\) and \(P_{\Pi,d}\) on the two evaluations satisfy
\[
 \|[P_{\Sigma,2},P_{\Pi,2}]\|=\frac{\sqrt2}{3},
 \qquad
 \|[P_{\Sigma,3},P_{\Pi,3}]\|=\frac{\sqrt3}{4}.
\]
In dimension three, the tritic mode of operative phase has singular value
\(s=1/2\), and
\[
 s^2=\frac14,\qquad 1-s^2=\frac34=\frac{90}{120},
 \qquad s\sqrt{1-s^2}=\frac{\sqrt3}{4}.
\]
Thus, the separation between the sheets is not added afterward: it is quantified in
the APP support itself.

\section{Discrete curvature sum--product}

In residues, write
\[
 S(x,y)=x+y,\qquad P(x,y)=xy\pmod9.
\]
For a potential \(T\), the progressive differences are defined
\[
 \Delta_xT(x,y)=T(x+1,y)-T(x,y),\qquad
 \Delta_yT(x,y)=T(x,y+1)-T(x,y).
\]
With \(A_x^T=\Delta_xT\) and \(A_y^T=\Delta_yT\), the curvature of an ordered
mixture is
\[
 F(T_x,T_y)=\Delta_xA_y^{T_y}-\Delta_yA_x^{T_x}.
\]
Since
\[
 \Delta_xS=\Delta_yS=1,\qquad
 \Delta_xP=y,\qquad \Delta_yP=x,
\]
it is obtained exactly
\begin{equation}
 F(S,S)=F(P,P)=0,\qquad
 F(S,P)=+1,\qquad F(P,S)=-1.
 \label{espiral:eq:curvatura-mixta-app}
\end{equation}

\begin{theorem}[Mixed oriented curvature of APP]
Pure compositions of the APP leaves are flat in the preceding
plaquette convention; mixed compositions have constant curvature of
opposite sign according to the order.
\end{theorem}
\begin{proof}
Substituting the four differences in the definition,
\[
 \Delta_x\Delta_yS-\Delta_y\Delta_xS=0,
 \quad
 \Delta_x\Delta_yP-\Delta_y\Delta_xP=0
\]
for the two pure compositions, while
\[
 \Delta_x\Delta_yP-\Delta_y\Delta_xS=1,
 \qquad
 \Delta_x\Delta_yS-\Delta_y\Delta_xP=-1.
\]
The identity  verifica in each to of the \(81\) plaquetas.
\end{proof}

For a coordinate rectangle with sides \(a,b\), cancellation of the
interior edges yields the finite form of Stokes' theorem:
\[
 W_{\partial R_{a,b}}(S,P)=ab,\qquad
 W_{\partial R_{a,b}}(P,S)=-ab\pmod9.
\]

\begin{figure}[htbp]
\centering
\begin{tikzpicture}[scale=.72]
\foreach \x in {0,...,8}{
  \draw[HMTLine] (\x,0)--(\x,8);
  \draw[HMTLine] (0,\x)--(8,\x);
}
\draw[HMTNavy,thick] (0,0) rectangle (8,8);
\draw[->,HMTBlue,very thick] (2,2)--(5,2) node[midway,below] {$S$};
\draw[->,HMTRed,very thick] (5,2)--(5,5) node[midway,right] {$P$};
\draw[->,HMTBlue!70!black,very thick] (5,5)--(2,5);
\draw[->,HMTRed!80!black,very thick] (2,5)--(2,2);
\node[fill=white,inner sep=2pt] at (3.5,3.5) {$F(S,P)=+1$};
\node[below,HMTGray,font=\small\sffamily] at (4,-.75) {soporte APP \(9\times9\)};
\end{tikzpicture}
\caption{Curvature arises from the order of the leaves over the same support.}
\label{espiral:fig:curvatura-app}
\end{figure}

\section{APP antecedents of radial classification}

APP contributes three distinct ingredients:
\begin{enumerate}
\item a domain bidimensional anterior to \(\R^2\);
\item two operations whose ordered mixture distinguishes orientation;
\item a ledger of residue, quotient, sheet, and carry capable of remembering the
path that produced the same output.
\end{enumerate}
The curvature \(\pm1\) is the local antecedent of the radial orientation, but
it must not be identified without proof with the global degree \(p\). The passage from one to
the other requires the TPK accumulation and the cocycle declared in the
\cref{espiral:chap:cociclos-radiales}.

\chapter{TRIT: local orientation and 3 geometric regimes}

\section{Lifted ternary state}

The visible observable is
\[
 \TRIT=\{-1,0,+1\}.
\]
In a local step with a carry, the balanced decomposition is written
\[
 \widehat\tau=(r,c),\qquad a+b=r+3c,
 \qquad r\in\{-1,0,+1\},\ c\in\Z.
\]
The visible zero possesses at least three elementary lifts:
\[
 0^+=(0,+1),\qquad0^0=(0,0),\qquad0^-=(0,-1).
\]
The visible equality does not fix the continuation. This also depends on sheet,
orientation, height, carry, boundary, path and phase.

\section{Quadratic algebras}

Para \(\tau\in\TRIT\), sea
\[
 \mathbb A_\tau=\R[J]/(J^2+\tau I).
\]
Three regimes are obtained:
\[
 J_{+1}^2=-I,\qquad J_0^2=0,\qquad J_{-1}^2=+I,
\]
and a single exponential identity
\[
 \exp(tJ_\tau)=C_\tau(t)I+S_\tau(t)J_\tau,
\]
where
\[
\begin{array}{c|c|c|c}
\tau & J_\tau^2 & C_\tau(t)&S_\tau(t)\\ \hline
+1&-I&\cos t&\sin t\\
0&0&1&t\\
-1&I&\cosh t&\sinh t
\end{array}
\]

\begin{figure}[htbp]
\centering
\begin{tikzpicture}[>=Latex]
\begin{scope}[xshift=-5cm]
 \draw[->,HMTGray] (-1.8,0)--(1.8,0);
 \draw[->,HMTGray] (0,-1.8)--(0,1.8);
 \draw[HMTBlue,very thick] (0,0) circle (1.25);
 \draw[->,HMTRed,thick] (0,0)--({1.25*cos(40)},{1.25*sin(40)});
 \node[below,HMTNavy,font=\sffamily\bfseries] at (0,-2.05) {elliptic \(J^2=-I\)};
\end{scope}
\begin{scope}
 \draw[->,HMTGray] (-1.8,0)--(1.8,0);
 \draw[->,HMTGray] (0,-1.8)--(0,1.8);
 \draw[HMTBlue,very thick] (-1.3,-.65)--(1.3,.65);
 \draw[->,HMTRed,thick] (0,0)--(1.1,.55);
 \node[below,HMTNavy,font=\sffamily\bfseries] at (0,-2.05) {parabolic \(J^2=0\)};
\end{scope}
\begin{scope}[xshift=5cm]
 \draw[->,HMTGray] (-1.8,0)--(1.8,0);
 \draw[->,HMTGray] (0,-1.8)--(0,1.8);
 \draw[HMTBlue,very thick,domain=-1.3:1.3,samples=80] plot ({cosh(\x)},{sinh(\x)});
 \draw[HMTBlue,very thick,domain=-1.3:1.3,samples=80] plot ({-cosh(\x)},{sinh(\x)});
 \draw[->,HMTRed,thick] (0,0)--(1.25,.75);
 \node[below,HMTNavy,font=\sffamily\bfseries] at (0,-2.05) {hyperbolic \(J^2=I\)};
\end{scope}
\end{tikzpicture}
\caption{The three quadratic regimes of the TRIT. They are transport types, not
names of specific curves.}
\label{espiral:fig:trit-regimenes}
\end{figure}

\section{Separation between TRIT regime, radial character and Frenet curvature}

The theory requires three levels that must not be flattened:
\begin{enumerate}
\item \(\tau\) tipa the regimen local of transport;
\item \(p\) types the accumulated radial character of a history;
\item \(\kappa_{\mathrm{F}}\) and \(\mathcal T_{\mathrm{F}}\) describe the curve already
published in a real space.
\end{enumerate}
A circumference has elliptic angular holonomy, but the present TRIT
\(0\) is a parabolic threshold. The circumference and the present are not the
same object. Likewise, a published logarithmic spiral may include
segments whose local transitions traverse the three TRIT types.

\begin{principle}[Separation of levels]
No identity among \(\tau\), \(p\) and \(\kappa_{\mathrm{F}}\) is admitted by
similarity of sign, name, or drawing. Every relation must be constructed by means of
a map with declared domain, codomain, and action.
\end{principle}

\section{Orientation Inversion}

For a word \(\boldsymbol\tau=(\tau_1,\ldots,\tau_n)\),
\[
 C_{\mathrm{rev}}(\boldsymbol\tau)
 =(-\tau_n,\ldots,-\tau_1),
 \qquad C_{\mathrm{rev}}^2=\id.
\]
The involution exchanges the orientations and preserves the threshold. In the
sector of reversible routes it will induce inversion of the affine holonomy, but
that conclusion requires the local representation to satisfy
\(h_{e^\dagger}=h_e^{-1}\).

\chapter{Enriched state, unique continuum and double projection}

\section{Inverse history system}

Let \(\mathsf S_n\) be the admissible states at depth \(n\) and
\[
 \pi_{n+1,n}:\mathsf S_{n+1}\longrightarrow\mathsf S_n
\]
the truncations. The space of complete histories is
\[
 X_\infty=\varprojlim_n\mathsf S_n.
\]
With the ultrametric
\[
 d(x,y)=3^{-m(x,y)},
\]
where \(m(x,y)\) is the first depth at which they differ, the fixed-prefix
cylinders are clopen. The structure of the continuum is obtained by jointly
constructing, from the same APP--TRIT--TPK state, its five consubstantial
constructions, which emerge correlatively, and by passing through terminality,
naturality, and the limit. This publication neither disaggregates them nor
redefines them by means of curves.

\section{Archimedean Evaluation via Compatible Cylinders}

Let \(\widehat{\mathbb Q}\) be the uniform completion of \(\mathbb Q\), constructed as
classes of rational Cauchy sequences. The calibrated construction associates
with each finite state an integer direction \(k_n(s)\) and the rational
semi-closed cylinder
\[
 I_n(s)=\{q\in\mathbb Q:k_n(s)3^{-n}\le q<(k_n(s)+1)3^{-n}\}.
\]
Satisfies simultaneously:
\begin{enumerate}[label=\textup{(A\arabic*)}]
\item \(\bigcup_{s\in\mathsf S_n}I_n(s)=\mathbb Q\) for every \(n\);
\item if \(t\in\mathsf S_{n+1}\) truncates to \(s\), then
      \(I_{n+1}(t)\subseteq I_n(s)\), and the children partition the parent under the half-closed boundary convention;
\item there exist a constant \(C>0\) and \(0<\lambda<1\) such that
\[
 \operatorname{diam}I_n(s)\le C\lambda^n,
 \qquad 0<\lambda<1;
\]
\item every chain of nested cylinders arises from a compatible history,
      with the two boundary expansions identified by the
      Archimedean relation.
\end{enumerate}

\begin{theorem}[Archimedean Closure of the Genealogical Continuum]
\label{thm:espiral-cierre-arquimediano-continuo-genealogico}
Each calibrated history determines a unique value
\[
 \{\operatorname{val}(x)\}
 =\bigcap_n\overline{I_n(x_n)}^{\widehat{\mathbb Q}}.
\]
The evaluation
\[
 \operatorname{val}:X_\infty^{\mathrm{cal}}\longrightarrow\widehat{\mathbb Q}
\]
is continuous, Hölder, and surjective. The quotient by equality of value
satisfies
\[
 X_\infty^{\mathrm{cal}}/{\sim_{\arq}}
 \cong\widehat{\mathbb Q}\cong\R,
\]
where the last identification is the usual Archimedean recognition of
the completion of \(\mathbb Q\), not an input to the generator.
\end{theorem}
\begin{proof}
Contraction of diameters and completeness of \(\widehat{\mathbb Q}\) reduce each
nested chain to a point. Cylinder compatibility proves continuity; the
geometric diameter bound gives the Hölder estimate. By (A1) and (A2),
each rational Cauchy class selects at every depth a cylinder containing
it;
(A4) provides a compatible history and uniquely resolves the
boundaries. Surjectivity follows. Two histories have the same
image exactly when they belong to the relation \(\sim_{\arq}\), so
the induced map on the quotient is a continuous open bijection
onto the final cylinder topology; it is therefore a homeomorphism.
\end{proof}

The real line is, therefore, a proved output of the evaluation and not a
presupposed domain. The history is not reduced to its value.

\section{Double projection}

From a single calibrated domain they are distinguished
\[
 \mathfrak R:X_\infty^{\mathrm{cal}}\longrightarrow\R,
 \qquad
 \mathfrak E:X_\infty^{\mathrm{cal}}\times\mathscr C_{\mathrm{exc}}
 \longrightarrow\mathbf{Exc}.
\]
The first contracts compatible cylinders; the second preserves boundary
incidence toward the exceptional corridor. Spiral theory is inserted into the
first branch, but preserves in its ledger the data required to compare
the second. Value is not identified with incidence, nor is a publication made to
generate the continuum retrospectively.

\begin{figure}[htbp]
\centering
\begin{tikzpicture}[>=Latex,node distance=12mm and 13mm,
 n/.style={draw=HMTNavy,fill=white,minimum height=9mm,text width=34mm,align=center,font=\small\sffamily}]
\node[n] (x) {$X_\infty^{\mathrm{cal}}$\\common enriched state};
\node[n,below left=of x] (r) {Archimedean publication\\value and curve in \(\R^d\)};
\node[n,below right=of x] (e) {exceptional publication\\incidence and automorphisms};
\draw[->,HMTBlue,very thick] (x)--node[left,font=\small] {$\mathfrak R$} (r);
\draw[->,HMTRed,very thick] (x)--node[right,font=\small] {$\mathfrak E_c$} (e);
\draw[HMTGold,dashed,thick] (r)--node[below,font=\scriptsize\sffamily]{invariantes comunes} (e);
\end{tikzpicture}
\caption{The two projections are correlative and subsequent to the same state.}
\label{espiral:fig:doble-proyeccion-espiral}
\end{figure}

\section{5 regions of \texorpdfstring{\(\pi\)}{pi} and multiplicity of genealogies}

The microfiber of \(\pi\) contains five regions distinct with
multiplicities
\[
 432+4\cdot144=1008.
\]
The same visible value conserves different regional histories. This fact
provides the exact precedent of the spiral classification: equality of
a published curve does not force equality of leaf, carry, boundary, or
memory. The theory does not invent that multiplicity; it transports it to a new
type of geometric output.

\begin{resultbox}[title={Causal antecedence of the discrete continuum with respect to the radial reader}]
Thus far, the domain, sheets, regimes,
holonomy, enriched state, and publication of the continuum have been constructed. Only now
is it legitimate to introduce real coordinates \((r,\theta)\) and recognize conventional
curves.
\end{resultbox}
\part{Genealogical theory of spiral correspondences}

\chapter{Affine universal envelope of the APP leaves}

\section{Construction of the universal radial module and its endomorphic envelope}

The operations \(u\mapsto u+a\) and \(u\mapsto\lambda u\) suggest the
affine group of the line. APP, however, contains noninvertible multiplicative
sectors and precedes real publication. The appropriate internal target
is therefore the monoid of affine endomorphisms of a universal radial module.

Let \(\widetilde X_{\APP}\) be the set of complete APP states, that is,
states that conserve both leaves, residue, quotient, and carry, and let us add an
absorbing symbol \(\bot\). We define the free module
\begin{equation}
 U_{\APP}:=\Z^{(\widetilde X_{\APP}\sqcup\{\bot\})}/\langle[\bot]\rangle.
 \label{espiral:eq:modulo-radial-universal}
\end{equation}
For each enriched transition \(e\), with source \(x_e\) and active digits
\((i_e,j_e)\), APP produces, without decimal evaluation, the records
\[
 i_e+j_e=R_e^++9q_e^+,
 \qquad
 i_ej_e=R_e^\times+9q_e^\times.
\]
Let \(s_e^\Sigma(x_e)\) and \(s_e^\Pi(x_e)\) be the complete states obtained
by adjoining, respectively, \((R_e^+,q_e^+)\) and
\((R_e^\times,q_e^\times)\), while preserving the other sheet in the ledger. The
additive lift determines the explicit vector
\[
 A_e=[s_e^\Sigma(x_e)]-[x_e]\in U_{\APP},
\]
and the multiplicative lifting determines the total map
\(m_e:\widetilde X_{\APP}\sqcup\{\bot\}\to
\widetilde X_{\APP}\sqcup\{\bot\}\), extended by
\[
 m_e(x)=
 \begin{cases}
 s_e^\Pi(x_e),&x=x_e,\\
 \bot,&x\neq x_e,
 \end{cases}
 \qquad m_e(\bot)=\bot.
\]
Its linearization is
\[
 M_e([x])=[m_e(x)]\in U_{\APP}.
\]
Thus, \(A_e\) and \(M_e\) are not parameters free: are the two publications
linearized of the same transition APP complete.

We define
\[
 \AffEnd(U_{\APP})
 =\operatorname{End}(U_{\APP})\ltimes U_{\APP}
\]
with composition
\begin{equation}
 (\Lambda_2,C_2)\circ(\Lambda_1,C_1)
 =\bigl(\Lambda_2\Lambda_1,
 C_2+\Lambda_2C_1\bigr).
 \label{espiral:eq:producto-affend}
\end{equation}
The pair acts by \(u\mapsto\Lambda u+C\). The linear part may be
noninvertible. The subgroupoid of units is obtained by restricting
\(\Lambda\in\operatorname{Aut}(U_{\APP})\).

\begin{proposition}[Universal property]
Every map from the set of complete states to a \(V\)-module that vanishes
on \(\bot\) extends uniquely to a morphism
\(\chi:U_{\APP}\to V\). If it also intertwines the multiplicative
lifts, \(\chi M_e=M_e^V\chi\), then it uniquely transports
the affine factors \(h_e\) and all their prefix holonomies.
\end{proposition}
\begin{proof}
The first statement is the universal property of the free module quotiented
by \([\bot]=0\). The second follows by applying \(\chi\) to the semidirect
composition: \(\chi(C_2+\Lambda_2C_1)=
\chi(C_2)+\Lambda_2^V\chi(C_1)\).
\end{proof}

\section{Local representation of a transition}

Let \(e:x\to y\) be a TPK transition. The tritic operational-phase selector
determines whether the additive leaf, the multiplicative leaf, or the threshold acts. It is defined
\begin{equation}
 h_e=
 \begin{cases}
 (I,A_e),&\text{sheet }\Sigma,\\
 (M_e,0),&\text{sheet }\Pi,\\
 (I,0),&\text{threshold},
 \end{cases}
 \label{espiral:eq:factor-afin-local}
\end{equation}
where \(A_e\) and \(M_e\) are the complete APP lifts constructed
above. Their residues alone are not taken: the quotient and carry remain
in the ledger.

\begin{definition}[Internal radial representation]
An internal radial representation is an assignment
\[
 \rho_{\rad}:\mathscr P_{\TPK}^{\enr}
 \longrightarrow\AffEnd(U_{\APP})
\]
which sends each arrow to \(h_e\), the identity to \((I,0)\) and composition
of paths to the product \cref{espiral:eq:producto-affend}.
\end{definition}

The definition does not consult a curve. The factors arise from the sheet and the
lift that the state already contains. A subsequent numerical realization
may send \(U_{\APP}\) to \(\R\), but that evaluation does not modify the
internal representation.

\section{Prefix holonomies}

For
\[
 \gamma_N=e_N\cdots e_1,
\]
we define
\[
 H_0=(I,0),\qquad
 H_j=h_{e_j}\cdots h_{e_1}=(\Lambda_j,C_j).
\]

\begin{theorem}[Affine holonomy functorial]
For composable paths \(\gamma_1,\gamma_2\),
\[
 H_{\gamma_2\gamma_1}=H_{\gamma_2}H_{\gamma_1}.
\]
If \(H_{\gamma_i}=(\Lambda_i,C_i)\), then
\begin{equation}
 \Lambda_{\gamma_2\gamma_1}=\Lambda_2\Lambda_1,
 \qquad
 C_{\gamma_2\gamma_1}=C_2+\Lambda_2C_1.
 \label{espiral:eq:holonomia-afin-functorial}
\end{equation}
\end{theorem}
\begin{proof}
The successive action on \(u\) is
\[
 u\xmapsto{H_{\gamma_1}}\Lambda_1u+C_1
 \xmapsto{H_{\gamma_2}}
 \Lambda_2\Lambda_1u+C_2+\Lambda_2C_1.
\]
Uniqueness of the affine decomposition gives \cref{espiral:eq:holonomia-afin-functorial}.
\end{proof}

For a nine-step block,
\begin{align}
 \Lambda_9&=\Lambda_{e_9}\cdots\Lambda_{e_1},\\
 C_9&=C_{e_9}+\Lambda_{e_9}C_{e_8}
 +\Lambda_{e_9}\Lambda_{e_8}C_{e_7}+\cdots
 +\Lambda_{e_9}\cdots\Lambda_{e_2}C_{e_1}.
 \label{espiral:eq:traslacion-nueve}
\end{align}
The position of each factor matters. This is the affine form of the memory of
order.

\section{Sum--product commutator}

In the invertible sector and on a one-dimensional realization, let
\[
 A_a(u)=u+a,\qquad M_\lambda(u)=\lambda u.
\]
Then
\begin{equation}
 A_aM_\lambda A_a^{-1}M_\lambda^{-1}
 =A_{a(1-\lambda)}.
 \label{espiral:eq:conmutador-afin}
\end{equation}
The proof is direct:
\[
 u\mapsto\lambda^{-1}u\mapsto\lambda^{-1}u-a
 \mapsto u-\lambda a\mapsto u+a(1-\lambda).
\]
The identity is asserted only when \(\lambda\) is a unit. In the full
monoid, the commutation defect is preserved by the factors and the
ledger, not by a commutator requiring nonexistent inverses.

\begin{figure}[!htb]
\centering
\begin{tikzpicture}[>=Latex,node distance=18mm and 28mm,
 n/.style={draw=HMTNavy,fill=white,minimum width=23mm,minimum height=8mm,font=\small}]
\node[n] (u) {$u$};
\node[n,right=of u,yshift=10mm] (a) {$u+a$};
\node[n,right=of a] (ma) {$\lambda u+\lambda a$};
\node[n,right=of u,yshift=-10mm] (m) {$\lambda u$};
\node[n,right=of m] (am) {$\lambda u+a$};
\draw[->,HMTBlue,thick] (u)--node[above left] {$A_a$} (a);
\draw[->,HMTRed,thick] (a)--node[above] {$M_\lambda$} (ma);
\draw[->,HMTRed,thick] (u)--node[below left] {$M_\lambda$} (m);
\draw[->,HMTBlue,thick] (m)--node[below] {$A_a$} (am);
\draw[->,HMTGold,dashed,thick] (ma)--node[right,font=\scriptsize]
  {$A_{a(1-\lambda)}$} (am);
\node[below=12mm of u,HMTGray,font=\small\sffamily,text width=80mm,align=center]
{The two orders end at distinct points; the vertical arrow measures
exactly the composition defect.};
\end{tikzpicture}
\caption{The sum-product difference manifests as affine holonomy.}
\label{espiral:fig:cuadrado-afin}
\end{figure}

\chapter{Radial cocycles in the TRIT lift of 2 semicoronas}
\label{espiral:chap:cociclos-radiales}

\section{Lifted category and elementary characters}

Let \(\mathscr P_{\TPK}^{2\mathrm c,\enr}\) be the category of paths whose
arrows are triples
\[
 \widehat e=(e,\tau_+(e),\tau_-(e)),
 \qquad \tau_\pm(e)\in\{-1,0,+1\},
\]
where the oriented pair belongs to the enriched state and the forgetful functor
sends \(\widehat e\) to the TPK transition \(e\). We define on each generator
\begin{equation}
 p(\widehat e)=\tau_+(e)+\tau_-(e),
 \qquad
 a(\widehat e)=\tau_+(e)-\tau_-(e).
 \label{espiral:eq:carta-doble-trit}
\end{equation}
For a route \(\widehat\gamma_N=\widehat e_N\cdots\widehat e_1\), the
accumulated cocycles are
\begin{equation}
 P_j=\sum_{k=1}^{j}p(\widehat e_k),
 \qquad
 A_j=\sum_{k=1}^{j}a(\widehat e_k),
 \qquad P_0=A_0=0.
 \label{espiral:eq:cociclos-radiales}
\end{equation}
Additivity under concatenation proves that \((P,A)\) is a \(\Z^2\)-valued cocycle of the lifted category. This assertion does not identify the pair by fiat with another previously named cocycle: it constructs the pair from the two orientations actually preserved along each arrow. Subdivision by stationary micro-arrows preserves both sums.

\section{Local exhaustive chart of the radial and oriented characters}

\begin{longtable}{cccccc}
\caption{Complete local chart of the two TRIT orientations.}
\label{espiral:tab:doble-trit}\\
\toprule
\(\tau_+\)&\(\tau_-\)&\(p\)&\(a\)&\(r_3(p)\)&\(c_3(p)\)\\
\midrule
\endfirsthead
\toprule
\(\tau_+\)&\(\tau_-\)&\(p\)&\(a\)&\(r_3(p)\)&\(c_3(p)\)\\
\midrule
\endhead
+1&+1&+2&0&-1&1\\
+1&0&+1&+1&+1&0\\
+1&-1&0&+2&0&0\\
0&+1&+1&-1&+1&0\\
0&0&0&0&0&0\\
0&-1&-1&+1&-1&0\\
-1&+1&0&-2&0&0\\
-1&0&-1&-1&-1&0\\
-1&-1&-2&0&+1&-1\\
\bottomrule
\end{longtable}

The three states with \(p=0\) are distinguished by
\(a=2,0,-2\). Therefore the radial degree is not sufficient to
reconstruct the orientation.

\section{Balanced decomposition}

For every \(p\in\Z\), there exists to unique decomposition
\[
 p=r_3(p)+3c_3(p),
 \qquad r_3(p)\in\{-1,0,+1\},\quad c_3(p)\in\Z.
\]
In the local block \(p\in[-2,2]\), the last column of
\cref{espiral:tab:doble-trit} gives the carry. In long routes, \(c_3(p)\) does not remain
restricted to \(\{-1,0,+1\}\). This precision avoids converting a local
example into a false global law.

\section{Involution}

Local Involution
\[
 (\tau_+,\tau_-)\longmapsto(-\tau_-,-\tau_+)
\]
yields
\[
 p\longmapsto-p,
 \qquad
 a\longmapsto a.
\]
The radial character changes orientation, whereas the oriented difference
remains. This separation will be visible in the identity
\(L_{-p}(x^{-1})=-L_p(x)\).

\chapter{Power-logarithmic coordinate and sum-product law}

\section{Definition and domain}

For \(x>0\), define
\begin{equation}
 L_p(x)=
 \begin{cases}
 \dfrac{x^p-1}{p},&p\neq0,\\[2mm]
 \log x,&p=0.
 \end{cases}
 \label{espiral:eq:lp-def}
\end{equation}
Its inverse is
\begin{equation}
 \exp_p(u)=
 \begin{cases}
 (1+pu)^{1/p},&p\neq0,\quad1+pu>0,\\[1mm]
 \e^u,&p=0.
 \end{cases}
 \label{espiral:eq:exp-p}
\end{equation}
The condition \(1+pu>0\) is not ornamental: it fixes the positive real branch.

\section{Deformed addition}

In the corresponding domain, let
\[
 u\boxplus_pv=u+v+puv.
\]

\begin{theorem}[Sum--Product Law]
For \(x,y>0\),
\begin{equation}
 L_p(xy)=L_p(x)\boxplus_pL_p(y).
 \label{espiral:eq:lp-producto}
\end{equation}
Moreover, \(0\) is the identity of \(\boxplus_p\), and
\[
 \exp_p(u\boxplus_pv)=\exp_p(u)\exp_p(v).
\]
\end{theorem}
\begin{proof}
Para \(p\neq0\),
\[
 1+pL_p(xy)=(xy)^p=x^py^p
 =(1+pL_p(x))(1+pL_p(y)).
\]
Expanding the last product and dividing by \(p\) yields
\cref{espiral:eq:lp-producto}. For \(p=0\) it is the identity
\(\log(xy)=\log x+\log y\). The equality for \(\exp_p\) follows by applying
the inverse.
\end{proof}

This law exactly expresses the APP articulation: the product in the
radial coordinate is published as a deformed sum in the coordinate
\(L_p\). The term \(puv\) preserves the interaction that would disappear in
an ordinary sum.

\section{Bidirectional Inversion}

\begin{proposition}[Inversion Identity]
For \(p\in\Z\) and \(x>0\),
\begin{equation}
 L_{-p}(x^{-1})=-L_p(x).
 \label{espiral:eq:lp-inversion}
\end{equation}
\end{proposition}
\begin{proof}
Si \(p\neq0\),
\[
 L_{-p}(x^{-1})
 =\frac{(x^{-1})^{-p}-1}{-p}
 =-\frac{x^p-1}{p}.
\]
For \(p=0\),
\(\log(x^{-1})=-\log x\).
\end{proof}

The involution simultaneously exchanges expansion/contraction and
\(p/-p\). It does not eliminate the memory of orientation, which remains in
\(a\).

\section{Logarithmic limit}

Para \(x>0\),
\[
 \lim_{p\to0}\frac{x^p-1}{p}=\log x.
\]
Indeed, \(x^p=\e^{p\log x}=1+p\log x+O(p^2)\). Therefore, the
logarithmic branch is not an added exception: it is the continuous limit of the family.
In the genealogical theory, however, \(p=0\) continues to preserve
\(a\) and the ledger; the functional limit does not identify its three
local origins of \cref{espiral:tab:doble-trit}.

\chapter{Power--logarithmic classification of the published curves}

\section{Uniform translation in the radial coordinate}

Suppose, after the publication of the continuum, that
\[
 L_p(r/r_0)=\beta m,
 \qquad
 \theta=\theta_0+\omega m,
 \qquad \omega\neq0.
\]
Eliminando \(m\),
\begin{equation}
 r_p(\theta)=r_0
 \left[1+p\frac{\beta}{\omega}
 (\theta-\theta_0)\right]^{1/p},
 \qquad p\neq0,
 \label{espiral:eq:familia-potencia}
\end{equation}
and
\begin{equation}
 r_0(\theta)=r_0
 \exp\!\left[\frac{\beta}{\omega}
 (\theta-\theta_0)\right].
 \label{espiral:eq:familia-log}
\end{equation}

\begin{theorem}[Publication of the five canonical classes]
After reoriginating \(\theta\) and rescaling \(r\), the specializations of
\cref{espiral:eq:familia-potencia,espiral:eq:familia-log} are:
\[
\begin{array}{c|c|c}
p&\text{normalized equation}&\text{conventional recognition}\\ \hline
2&r^2=a^2\theta&\text{Fermat spiral}\\
1&r=a\theta&\text{Archimedean spiral}\\
0&r=r_0\e^{b\theta}&\text{logarithmic spiral}\\
-1&r=a/\theta&\text{hyperbolic spiral}\\
-2&r^2=a^2/\theta&\text{lituus}
\end{array}
\]
in the corresponding positive real domains.
\end{theorem}
\begin{proof}
For \(p=1\), \(r\) is affine in \(\theta\), and recentering removes the
constant term. For \(p=2\), \(r^2\) is affine. For \(p=-1\),
\(r^{-1}\) is affine; for \(p=-2\), \(r^{-2}\) is affine. The \(p=0\) branch
is \cref{espiral:eq:familia-log}. Positive constants are absorbed into \(a,b\)
without altering the exponent.
\end{proof}

\begin{figure}[!htb]
\centering
\begin{tikzpicture}[scale=.62]
\begin{scope}[xshift=0cm,yshift=0cm]
  \draw[->,HMTLine] (-1.9,0)--(1.9,0);
  \draw[->,HMTLine] (0,-1.9)--(0,1.9);
  \draw[HMTBlue,very thick,domain=.15:18.2,samples=180,smooth,variable=\t]
    plot ({sqrt(.14*(\t+1))*cos(deg(\t))},{sqrt(.14*(\t+1))*sin(deg(\t))});
  \node[below,text=HMTGray,font=\scriptsize\sffamily] at (0,-2.1) {Fermat, \(p=2\)};
\end{scope}
\begin{scope}[xshift=5cm,yshift=0cm]
  \draw[->,HMTLine] (-1.9,0)--(1.9,0);
  \draw[->,HMTLine] (0,-1.9)--(0,1.9);
  \draw[HMTNavy,very thick,domain=.15:18.2,samples=180,smooth,variable=\t]
    plot ({.08*(\t+1)*cos(deg(\t))},{.08*(\t+1)*sin(deg(\t))});
  \node[below,text=HMTGray,font=\scriptsize\sffamily] at (0,-2.1) {Archimedes, \(p=1\)};
\end{scope}
\begin{scope}[xshift=10cm,yshift=0cm]
  \draw[->,HMTLine] (-1.9,0)--(1.9,0);
  \draw[->,HMTLine] (0,-1.9)--(0,1.9);
  \draw[HMTRed,very thick,domain=.15:18.2,samples=180,smooth,variable=\t]
    plot ({.18*exp(.055*\t)*cos(deg(\t))},{.18*exp(.055*\t)*sin(deg(\t))});
  \node[below,text=HMTGray,font=\scriptsize\sffamily] at (0,-2.1) {Logarithmic, \(p=0\)};
\end{scope}
\begin{scope}[xshift=2.5cm,yshift=-5.2cm]
  \draw[->,HMTLine] (-1.9,0)--(1.9,0);
  \draw[->,HMTLine] (0,-1.9)--(0,1.9);
  \draw[HMTGold,very thick,domain=.15:18.2,samples=180,smooth,variable=\t]
    plot ({1.8/(\t+1.8)*cos(deg(\t))},{1.8/(\t+1.8)*sin(deg(\t))});
  \node[below,text=HMTGray,font=\scriptsize\sffamily] at (0,-2.1) {Hyperbolic, \(p=-1\)};
\end{scope}
\begin{scope}[xshift=7.5cm,yshift=-5.2cm]
  \draw[->,HMTLine] (-1.9,0)--(1.9,0);
  \draw[->,HMTLine] (0,-1.9)--(0,1.9);
  \draw[HMTGray,very thick,domain=.15:18.2,samples=180,smooth,variable=\t]
    plot ({sqrt(2.0/(\t+1.8))*cos(deg(\t))},{sqrt(2.0/(\t+1.8))*sin(deg(\t))});
  \node[below,text=HMTGray,font=\scriptsize\sffamily] at (0,-2.1) {Lituus, \(p=-2\)};
\end{scope}
\end{tikzpicture}
\caption{Five power--logarithmic family publications on common axes and scales.}
\label{espiral:fig:cinco-espirales}
\end{figure}

\section{Non-exhaustiveness and exact scope}

The theorem classifies the family generated by uniform translation in the
coordinate \(L_p\). It does not classify all possible plane curves. The
complete genealogical taxonomy must include at least
\[
 (p,a,[\Lambda,C],[\kappa_9],
 \omega_{\mathrm{ang}},\mu_{\mem},\varepsilon_{\mathrm{tw}}).
\]
Two routes with the same \(p\) may differ in self-scaling, orientation,
carry, or boundary and publish distinct curves. Two routes with all visible
data equal may still differ in an eliminated ledger coordinate.

\section{General affine holonomy}

If the coordinate \(u=L_p(r/r_0)\) satisfies
\[
 u_{n+1}=\Lambda u_n+C,
\]
the iteration internal exact, for \(n\in\N\), is
\[
 u_n=\Lambda^n u_0+\sum_{k=0}^{n-1}\Lambda^kC.
\]
If \(I-\Lambda\) is invertible, may iscribirse
\[
 u_n=u_*+\Lambda^n(u_0-u_*),
 \qquad u_*=(I-\Lambda)^{-1}C.
\]
Continuous interpolation is performed only after fixing a positive
scalar character \(\chi:U_{\APP}\to\R\) such that
\(\chi\Lambda=\lambda\chi\), with \(\lambda>0\). Writing
\(u=\chi(u)\), \(C=\chi(C)\), if \(\lambda\neq1\), then
\(u_*=C/(1-\lambda)\). For \(\theta_n=\theta_0+n\omega\),
\(\omega\neq0\), the scalar realization admits
\begin{equation}
 r(\theta)=r_0\exp_p\!\left[
 u_*+\lambda^{(\theta-\theta_0)/\omega}(u_0-u_*)
 \right].
 \label{espiral:eq:familia-afin-general}
\end{equation}
Real powers belong exclusively to this realization; the
internal dynamics uses integer powers of \(\Lambda\). The pair
\((p,[\Lambda,C])\) enlarges the classical list: \(p\) fixes the radial chart, and
the affine conjugacy class fixes the dynamics within it.

\begin{remark}[Fixed point and electron]
The existence of an affine fixed point does not by itself identify that point with the
central electron \((5,5)\). The electron is a prior material result; an
identification with a radial fixed point requires a separate realization
square.
\end{remark}

\chapter{Radial enriched reader}

\section{Obstruction to the reduced reader}

The initially suggested map
\[
 X_{\TPK}^{\enr}\to(p,a,\Lambda,C,m,r,\theta)
\]
mixes causal levels and loses prefixs. The coordinates \(r,\theta\) are
arquimedianas, while \((\Lambda,C)\) preserves only the final product.

\begin{proposition}[The terminal pair admits no natural truncation]
In general there is no application that recovers the first factor of an
affine word from its final product.
\end{proposition}
\begin{proof}
In notation \((\Lambda,C)\),
\[
 (1,2)(2,0)=(2,2),
 \qquad
 (1,1)(2,1)=(2,2).
\]
The two final products coincide, but their initial prefixes are
\((2,0)\) and \((2,1)\). A function of the final product would have
to assign it two different truncations, which is impossible.
\end{proof}

\section{Definition of the complete reader}

Let
\(\widehat\gamma_N=\widehat e_N\cdots\widehat e_1\) one route of the
category lifted and let \((H_0,\Led_0)=((I,0),\Led(x_0))\) its datum
initial. We define
\begin{equation}
 \boxed{
 \mathcal P_{\rad,N}^{\enr}(\widehat\gamma_N)
 =\left(
 (H_0,\Led_0);
 \bigl(p(\widehat e_j),a(\widehat e_j),P_j,A_j,
 h_{e_j},H_j,\Led_j(\widehat\gamma_N)\bigr)_{j=1}^{N}
 \right).}
 \label{espiral:eq:lector-radial-enriquecido}
\end{equation}
The codomain \(\mathcal D_{\rad,N}^{\enr}\) contains compatible sequences
of those data, including the underlying edge stored in each entry of the
ledger and its source and target maps. The output conserves the initial datum,
each factor, each prefix holonomy, the local and accumulated cocycles, and every
memory update.

\begin{definition}[Admissible homogeneous radial sector]
Let \(\mathcal D_{\rad,\infty}^{\mathrm{hom}}\) the space of data marked
\[
 D=(D^{\enr};r_0,p,u_0,\beta,\theta_0,\omega),
\]
where \(D^{\enr}\in\mathcal D_{\rad,\infty}^{\enr}\), \(r_0>0\),
\(p\in\{-2,-1,0,1,2\}\), \(u_0,\beta,\theta_0,\omega\in\R\) and
\((\beta,\omega)\neq(0,0)\), talis that
\[
 p(\widehat e_j)=p,\qquad
 u(m)=u_0+\beta m,\qquad
 \theta(m)=\theta_0+\omega m.
\]
The mark is part of the data and is not chosen afterward; consequently, the evaluation
is a well-defined function.
Its real publication interval is
\[
 I_D=
 \begin{cases}
 \{m\in\R:1+p(u_0+\beta m)>0\},&p\neq0,\\
 \R,&p=0.
 \end{cases}
\]
Histories with variable degree remain in the enriched ledger and are published through an atlas of piecewise charts; they do not determine a single curve of constant degree.
\end{definition}

\begin{definition}[Homogeneous spiral publication]
After constructing the continuum and choosing the radial and angular characters compatible, it is defined
\[
 \ev_{\arq}^{\mathrm{esp}}:
 \mathcal D_{\rad,\infty}^{\mathrm{hom}}\longrightarrow
 \coprod_{D}C^\infty(I_D,\R^2)
\]
by
\[
 D\longmapsto
 \left[m\longmapsto
 \bigl(r_0\exp_{p}(u_0+\beta m)\cos(\theta_0+\omega m),
 r_0\exp_{p}(u_0+\beta m)\sin(\theta_0+\omega m)\bigr)\right].
\]
\end{definition}

The internal reader and external evaluation are separate. The first arrow preserves genealogy; the second may forget part of it. The condition \((\beta,\omega)\neq(0,0)\) excludes only the degenerate constant parametrization; on \(I_D\) the published curve is regular.

\section{Conservation as traceability}

The projection onto carry, boundary, and memory satisfies
\[
 \operatorname{pr}_{\mathrm{carry},frontier,memory}
 \circ\mathcal P_{\rad,N}^{\enr}
 =\Led_{{\mathrm{carry},frontier,memory},N}.
\]
To preserve does not mean to make those coordinates constant. It means that:
\[
 g(K+9)=g(K),
\]
\[
 q_{\mem}(\Gamma_9x)=q_{\mem}(x)+1,
\]
\[
 B_{K+9}=\operatorname{Tra}_{\Gamma_9}(B_K),
\]
and the carry is updated by \cref{espiral:eq:cociclo-carry}. No update
is hidden by the final curve.

\chapter{Naturalness, equivariance, and inverse limit}

\section{Naturalness under Truncation}

For \(M\le N\), let
\[
 \pi_{N,M}:\Hist_N^{\enr}(\TPK)\to\Hist_M^{\enr}(\TPK)
\]
the truncation of histories, and let
\[
 q_{N,M}:\mathcal D_{\rad,N}^{\enr}\to
 \mathcal D_{\rad,M}^{\enr}
\]
the removal of data with index greater than \(M\).

\begin{theorem}[Radial Reader Naturalness]
The diagram
\[
\begin{array}{ccc}
\Hist_N^{\enr}(\TPK)
 & \xrightarrow{\ \mathcal P_{\rad,N}^{\enr}\ }
 & \mathcal D_{\rad,N}^{\enr}\\[1.0em]
{\scriptstyle \pi_{N,M}}\big\downarrow
 && \big\downarrow{\scriptstyle q_{N,M}}\\[1.0em]
\Hist_M^{\enr}(\TPK)
 & \xrightarrow{\ \mathcal P_{\rad,M}^{\enr}\ }
 & \mathcal D_{\rad,M}^{\enr}
\end{array}
\]
commutes; that is,
\begin{equation}
 q_{N,M}\circ\mathcal P_{\rad,N}^{\enr}
 =\mathcal P_{\rad,M}^{\enr}\circ\pi_{N,M}.
 \label{espiral:eq:naturalidad-truncamiento}
\end{equation}
\end{theorem}
\begin{proof}
Both sides preserve exactly the initial record and the first
\(M\) arrows, together with their local and cumulative cocycles, affine factors,
prefix holonomies, and ledger entries of \(\widehat\gamma_N\). No
reconstruction from the final product is involved: the complete sequence
is part of the definition.
\end{proof}

\section{Reader of the inverse limit}

The compatibility \cref{espiral:eq:naturalidad-truncamiento} induces
\begin{equation}
 \mathcal P_{\rad,\infty}^{\enr}:
 \varprojlim_N\Hist_N^{\enr}(\TPK)
 \longrightarrow
 \varprojlim_N\mathcal D_{\rad,N}^{\enr}.
 \label{espiral:eq:lector-limite-inverso}
\end{equation}
Each finite component is obtained by truncation of the complete history.
Therefore the reader does not invent an external continuity: it prolongs the finite
certificates compatibly.

\section{Quadruple subdivision}

In the refinement \(\sd_4\) of the TPK, three micro-arrows are stationary and the
fourth reproduces the observable transition. If
\[
 \rho_{\rad}(e_1)=\rho_{\rad}(e_2)=\rho_{\rad}(e_3)=I,
 \qquad \rho_{\rad}(e_4)=h_e,
\]
then
\[
 H_{\sd_4(e)}=I\,I\,I\,h_e=h_e.
\]
Naturality is thus certified for this refinement. For a general substitution \(e\mapsto\widetilde e_k\cdots\widetilde e_1\), the exact condition is
\[
 h_e=h_{\widetilde e_k}\cdots h_{\widetilde e_1}.
\]
This equality, rather than the similarity of the published curves, determines the
compatibility of a new refinement.

\section{Equivariance with respect to nonadic holonomy}

Let \(b_9(x)\) be the block of nine transitions that realizes \(\Gamma_9\) from
\(x\), and
\[
 H_9(x)=H_{b_9(x)}=(\Lambda_9(x),C_9(x)).
\]
The action induced on the radial ledger adds that block.

\begin{theorem}[Cocycle Equivariance]
One has
\begin{equation}
 \mathcal P_{\rad}^{\enr}\circ\Gamma_9
 =\widehat\Gamma_9^{\rad}\circ
 \mathcal P_{\rad}^{\enr}.
 \label{espiral:eq:equivarianza-gamma9}
\end{equation}
In components,
\[
 P\mapsto P+p(b_9),
 \qquad
 A\mapsto A+a(b_9),
\]
\[
 \Lambda\mapsto\Lambda_9(x)\Lambda,
 \qquad
 C\mapsto C_9(x)+\Lambda_9(x)C,
\]
\[
 g\mapsto g,\qquad q_{\mem}\mapsto q_{\mem}+1.
\]
\end{theorem}
\begin{proof}
The action of \(\Gamma_9\) concatenates \(b_9(x)\) to the history, where
\(p(b_9)\) and \(a(b_9)\) are the sums of the local characters of its nine
lifted arrows. The cocycle law adds the degrees, and the affine law
\cref{espiral:eq:producto-affend} composes the holonomies. The TPK supplies the phase
return and the memory increment. The definition of
\(\widehat\Gamma_9^{\rad}\) reproduces precisely these updates.
\end{proof}

This is not literal commutation of endomorphisms of the same space. The
codomains are distinct and the memory changes. The correct term is
\emph{equivariance}.

\section{Investment in the reversible sector}

If \(h_{e^\dagger}=h_e^{-1}\), the inversion diagram satisfies
\[
 \mathcal P_{\rad}^{\enr}(C_{\mathrm{rev}}\gamma)
 =\mathcal J_{\rad}
 \mathcal P_{\rad}^{\enr}(\gamma),
\]
where \(\mathcal J_{\rad}\) reverses the order of factors, replaces each
factor by its inverse, and transports the ledger through the TPK involution. In
the real output, \cref{espiral:eq:lp-inversion} exchanges the two orientations.

\chapter{Genealogical classification and non-faithfulness of publication}

\section{Full radial signature of a history}

For a finite history
\(\widehat\gamma_N=\widehat e_N\cdots\widehat e_1\), define
\begin{equation}
 \mathfrak I_N(\widehat\gamma_N)=
 \left(
 (x_0,H_0,\Led_0);
 \bigl(\widehat e_j,p(\widehat e_j),a(\widehat e_j),P_j,A_j,
 h_{e_j},H_j,\Led_j\bigr)_{j=1}^{N}
 \right).
 \label{espiral:eq:invariante-espiral}
\end{equation}
The signature is complete relative to the chosen enriched radial reader: it preserves
the initial object, the typed sequence of arrows, the local and
accumulated characters, all affine prefixes, and the complete ledger. It is not presented
as a complete invariant of all HMT, nor does it replace the coordinates of the
enriched state that do not lie in the reader's domain.

\begin{definition}[Spiral Genealogical Equivalence]
Two histories are equivalent if there is a typed isomorphism of their
radial data that transports the sequence of
\cref{espiral:eq:invariante-espiral} term by term, respects sources, targets, composition, and
ledgers, and commutes with all declared truncations.
\end{definition}

The inverse limit yields the compatible family
\((\mathfrak I_N)_{N\ge0}\). Equality of curves in \(\R^2\) is a
coarser relation: the Archimedean evaluation factors through genealogical
equivalence, but in general does not reflect it.

\section{Classification by 3 visible increments}

Let \(\omega\) be the angular increment, \(\beta\) the radial increment in
the coordinate \(L_p\), and \(\mu\) the lifted-memory increment. The minimal
visible table is:
\begin{longtable}{ccccp{.34\textwidth}}
\caption{Kinematic classification of publications.}
\label{espiral:tab:clasificacion-cinematica}\\
\toprule
\(\omega\)&\(\beta\)&\(\mu\)&Published type&Interpretation\\
\midrule
\endfirsthead
\toprule
\(\omega\)&\(\beta\)&\(\mu\)&Published type&Interpretation\\
\midrule
\endhead
0&\(\neq0\)&0&recta radial&scale variation without rotation or memory\\
\(\neq0\)&0&0&circle&intrinsic radial angular closure\\
\(\neq0\)&0&\(\neq0\)&cylindrical helix&periodic phase with raised memory\\
\(\neq0\)&\(\neq0\)&0&espiral plana&radial variation without visible memory axis\\
\(\neq0\)&\(\neq0\)&\(\neq0\)&radial helical suspension&turning, self-scaling and memory\\
0&0&\(\neq0\)&fibra axial&memory advance without transverse figure\\
\bottomrule
\end{longtable}

This table classifies the published kinematics, not the entire genealogy.
Two rows with the same three increments may differ in \(p\), carry,
boundary, or affine class.

\section{3 fibers of the same visible evaluation}

\begin{example}[Circle with null memory]
Let \(H_j=I\), \(\mu=0\), and \(\theta_j=2\pi j/9\) be given. The output is a
circumference, and the ledger does not register axial lift.
\end{example}

\begin{example}[The same circumference as shadow]
Let \(H_j=I\), \(\mu=1\) and the same angular phase. En \(\R^3\) is obtained the
helix \cref{espiral:eq:inmersion-helice-nonadica}; to the aplicar \(\pi_{xy}\), the
output flat coincide with the of the example preceding.
\end{example}

\begin{example}[Same spiral, different carry]
Let there be two histories with the same real characters \(p,\beta,\omega\)
but with different carry decompositions that radial evaluation maps
to the same increment \(\beta\). Both publish \(r_p(\theta)\), while their
ledgers differ. The enriched reader separates them; the curve does not.
\end{example}

\section{Genealogical Classification Theorem}

\begin{theorem}[Spiral Genealogical Correspondences]
Once the universal radial module, the local representation
\(e\mapsto h_e\), and the radial cocycles of the TRIT lift are fixed, every
compatible TPK history determines:
\begin{enumerate}
\item one natural sequence of factors and prefix affine holonomies;
\item to class genealogical istable under truncamiento;
\item an equivariant action of the nonadic holonomy;
\item in the reversible sector, one conjugate history of inverse holonomy;
\item for each admissible homogeneous radial sector, a regular curve in
\(\R^2\) and, by conserving the depth, a suspension in \(\R^3\).
\end{enumerate}
A history of variable degree instead determines an ordered atlas of
piecewise charts. The application to curves is not faithful: it may identify
histories with distinct ledgers.
\end{theorem}
\begin{proof}
Items 1 and 2 are the affine-holonomy and naturality theorems. Item
3 is \cref{espiral:eq:equivarianza-gamma9}; item 4 is inversion in the
reversible subgroupoid. Item 5 is obtained by restricting the inverse
limit to the domain \(\mathcal D_{\rad,\infty}^{\mathrm{hom}}\) and applying its
radial and angular characters. Non-faithfulness follows from
\cref{espiral:fig:circulo-helice} and from the general proposition of the first chapter.
\end{proof}

\begin{statusbox}[title={Hypotheses, domain, and scope of the theorem}]
The nonadic holonomy, the helix, the Pell screw, and the TPK cocycles
are recovered results. The affine envelope, the prefix reader, its
naturality, and the publication \(L_p\) constitute a new exact
formalization in the declared domains. The local chart
\(p(\widehat e)=\tau_+(e)+\tau_-(e)\) and
\(a(\widehat e)=\tau_+(e)-\tau_-(e)\) is exact in the lifted category
of 2 semicrowns; its use outside that domain requires exhibiting the functor that
supplies both orientations. Publication as a single curve is restricted
to the admissible homogeneous sector defined in this chapter.
\end{statusbox}
